\documentclass[11pt]{article}
\usepackage[a4paper,margin=18mm]{geometry}
\usepackage{fontspec}
\setmainfont{Latin Modern Roman}
\usepackage{amsmath,amssymb,amsthm,amscd,graphicx,color}
\usepackage{xurl}
\usepackage[unicode,hidelinks]{hyperref}
\allowdisplaybreaks
\setlength\emergencystretch{3em}
\usepackage{enumerate,mathrsfs}
\usepackage[all,2cell]{xy}

\newtheorem{innercustomthm}{Theorem}
\def\drawline(#1,#2)(#3,#4){\qbezier(#1,#2)(#1,#2)(#3,#4)}


\makeatletter
%
\newtheorem{theorem}{Theorem}
\newtheorem*{theorem*}{Theorem}
\newtheorem{lemma}{Lemma}
\newtheorem*{lemma*}{Lemma}
\newtheorem{corollary}{Corollary}
\newtheorem*{corollary*}{Corollary}
\newtheorem{proposition}{Proposition}
\newtheorem*{proposition*}{Proposition}
\newtheorem{conjecture}{Conjecture}
\newtheorem*{conjecture*}{Conjecture}
{\theoremstyle{definition} \newtheorem{definition}{Definition}
\newtheorem*{definition*}{Definition}
\newtheorem{example}{Example}
\newtheorem*{example*}{Example}
\newtheorem{remark}{Remark}
\newtheorem*{remark*}{Remark}
\newtheorem{note}{Note}
\newtheorem*{note*}{Note}
\makeatother

\makeatletter
\newcommand{\dash}[1]{\text{-}\mathrm{#1}}

\newenvironment{customthm}[1]{\renewcommand\theinnercustomthm{#1}\innercustomthm}{\endinnercustomthm}
\makeatother
\begin{document}
\begin{center}\Large Extension fullness of finitely generated locally Gelfand-Zeitlin-finite gl\_n(C) modules in all U(gl\_n(C)) modules\end{center}
\noindent\textbf{Stable ID:} 2004\par
\noindent\textbf{Authors:} Kevin Coulembier; Volodymyr Mazorchuk\par
\noindent\textbf{Original work:} Extension Fullness of the Categories of Gelfand-Zeitlin and Whittaker Modules\par
\noindent\textbf{Version:} Official SIGMA 11 (2015) article 016, DOI 10.3842/SIGMA.2015.016; journal PDF and native TeX acquired 2026-10-01, exact bytes pinned by SHA256\par
\noindent\textbf{Selected task scope:} Full TheoremA/Section3.3 Theorem1: for n>=\allowbreak{}1 over C, the inclusion of finitely generated Gamma-locally-finite Gelfand-Zeitlin modules into all U(gl\_n)-modules induces bijections on Yoneda Ext\textasciicircum{}d for every d>=\allowbreak{}0 and every pair of objects. No assertion for arbitrary non-finitely-generated Gamma-locally-finite modules.\par
\noindent\textbf{Difficulty:} H2: The proof compares Yoneda extensions in categories lacking the usual supply of projective objects, using explicit pushout and pullback constructions for exact adjunction. It then combines generalized character decomposition, freeness over the Gelfand-Zeitlin algebra and a Serre-subcategory induction to prove all-degree Ext comparison for arbitrary pairs in the stated category.\par
\noindent\textbf{Editorial boundary note (not author text):} The complete original Theorem A/Section 3.3 Theorem 1 and full Gelfand–Zeitlin author proof are retained for finitely generated Gamma-locally-finite U(gl\_n(C))-modules. The stronger Serre-subcategory induction and entire exact-adjunction Yoneda pushout/pullback proof are supplied without presuming enough projective/injective objects. All Gamma/Harish-Chandra/freeness, generalized-character and induction/restriction definitions are given, while their named previously established properties remain cited prerequisites. Adjunction Proposition 4 keeps its original number; unneeded Propositions 2–3, Whittaker theorem and global-dimension corollaries are not selected or counted.\par
\noindent\textbf{Source:} \url{https://sigma-journal.com/2015/016/}\quad DOI: \url{https://doi.org/10.3842/SIGMA.2015.016}\par
\noindent\textbf{License and attribution:} Copyright retained by the named authors. Published by SIGMA under CC BY 4.0: \url{https://creativecommons.org/licenses/by/4.0/}. Source: \url{https://sigma-journal.com/about.html}. Adaptation consists of standalone excerpt formatting, cover and numbering restoration; original author language and mathematical text are retained.\par
\noindent\textbf{Editorial symbol notice (not author text):} In the exact-adjunction proof a prose expression C(N):=\allowbreak{}GF(M)/N-prime has an M/N variable slip; the displayed diagram and the definition of C as the cokernel of Id\_A->GF specify the objects. No source formula is altered.\par
\noindent\textbf{Typesetting adaptation (not author text):} The unavailable epic package is replaced only for the original straight-line embedding diagram by standard picture qbezier strokes on the same author coordinates. Diagram text and mathematical proof are unchanged.\par
\medskip\hrule\medskip\noindent\textbf{Author text begins below.}\par
\setcounter{section}{1}
\setcounter{subsection}{0}
\setcounter{subsubsection}{0}
\setcounter{innercustomthm}{0}
\setcounter{equation}{0}
% Original source lines 109--123
The f\/irst main result of this paper is the following statement proved in Section~\ref{s2} (we refer to Sections~\ref{s1}
and~\ref{s2} for more details):

\begin{customthm}{A}
\label{tmain1}
The category $\mathscr{G}\mathscr{Z}$ is extension full in the category of all $\mathfrak{gl}_n$-modules.
\end{customthm}

The added dif\/f\/iculty of the category $\mathscr{G}\mathscr{Z}$ in comparison with thick category $\mathcal{O}$
in~\cite{CoMa} is that~$\mathscr{G}\mathscr{Z}$ is not a~Serre subcategory generated by a~well-known category which has
enough projective objects (contrary to the relation between category~$\mathcal{O}$ and thick category $\mathcal{O}$).
Therefore to prove Theorem~\ref{tmain1} we have to modify and strengthen the abstract results on extension fullness
in~\cite{CoMa}.
Our arguments also heavily use some properties of Gelfand--Zeitlin modules established by Futorny and Ovsienko
in~\cite{FuOv2}.
\setcounter{section}{1}
\setcounter{subsection}{0}
\setcounter{subsubsection}{0}
\setcounter{innercustomthm}{0}
\setcounter{equation}{0}
% Original source lines 165--289
\section{Extension full subcategories}\label{s1}

\subsection{Extensions in Abelian categories}%\label{s1.1}

Let $\mathcal{A}$ be an Abelian category and $M,N\in\mathcal{A}$.
Recall that, for $d\in\mathbb{Z}_{\geq 0}$, the set $\mathrm{Ext}_{\mathcal{A}}^d(M,N)$ of {\em degree~$d$ extensions}
from~$M$ to~$N$ is def\/ined as the set of equivalence classes of exact sequences
\begin{gather*}
\xymatrix{
0\ar[r]&N\ar[r]&X_1\ar[r]&X_2\ar[r]&\cdots\ar[r]&X_d\ar[r]&M\ar[r]&0,
}
\end{gather*}
where all objects and all morphisms are in $\mathcal{A}$, modulo the minimal equivalence relation which contains the
binary relation given by existence of a~commutative diagram
\begin{gather*}
\xymatrix{
0\ar[r]&N\ar[r]\ar@{=}[d]&X_1\ar[r]\ar[d]&X_2\ar[r]\ar[d]&\cdots\ar[r]&X_d\ar[r]\ar[d]&M\ar[r]\ar@{=}[d]&0\\
0\ar[r]&N\ar[r]&Y_1\ar[r]&Y_2\ar[r]&\cdots\ar[r]&Y_d\ar[r]&M\ar[r]&0
}
\end{gather*}
The set $\mathrm{Ext}_{\mathcal{A}}^d(M,N)$ has the natural structure of an Abelian group via the Baer sum.
If $\mathcal{A}$ is $\Bbbk$-linear for some f\/ield $\Bbbk$, then $\mathrm{Ext}_{\mathcal{A}}^d(M,N)$ has the structure of
a~$\Bbbk$-vector space.
We refer to~\cite[Section~3.4]{We} for further information and details.

For $M\in\mathcal{A}$, the {\em projective dimension} $\mathrm{proj.dim}(M)$ of~$M$ is def\/ined as the maximal~$d$ such
that there is $N\in \mathcal{A}$ with $\mathrm{Ext}_{\mathcal{A}}^d(M,N)\neq 0$.
If such maximal~$d$ does not exist, then $\mathrm{proj.dim}(M):=\infty$.
Dually one def\/ines the {\em injective dimension} $\mathrm{inj.dim}(N)$ for $N\in\mathcal{A}$.
The {\em global dimension} $\mathrm{gl.dim}(\mathcal{A})\in\mathbb{Z}_{\geq 0}\cup\{\infty\}$ is the supremum of
projective dimensions taken over all objects in $\mathcal{A}$.
The global dimension coincides with the supremum of injective dimensions taken over all objects in $\mathcal{A}$
(see~\cite[Lemma~5.11.11]{Zi}).

\subsection{Extension full subcategories}\label{s1.2}

Let $\mathcal{A}$ be an Abelian category and $\mathcal{B}$ a~full Abelian subcategory of $\mathcal{A}$ in the sense that
the Abelian structure of $\mathcal{B}$ is inherited from $\mathcal{A}$.
In particular, the natural inclusion functor $\boldsymbol{\iota}:\mathcal{B}\to \mathcal{A}$ is exact.
Then, for every $M,N\in\mathcal{B}$ and every $d\in\mathbb{Z}_{\geq0}$, the functor $\boldsymbol{\iota}$ induces
homomorphisms
\begin{gather*}
\iota^d_{M,N}: \ \mathrm{Ext}_{\mathcal{B}}^d(M,N)\to \mathrm{Ext}_{\mathcal{A}}^d(M,N)
\end{gather*}
of Abelian groups.
In general, these homomorphisms $\iota^d_{M,N}$ are neither injective nor surjective.

We say that $\mathcal{B}$ is {\em extension full} in $\mathcal{A}$ provided that $\iota^d_{M,N}$ are bijective for all
$d\in\mathbb{Z}_{\geq0}$ and for all $M,N\in\mathcal{B}$.
We refer to~\cite[Section~2]{CoMa} for details.

Let $0\to K\to M\to N\to 0$ be a~short exact sequence in $\mathcal{B}$.
Then, for $Q\in \mathcal{B}$, application of $\mathrm{Hom}_{\mathcal{B}}(Q,{}_-)$ and
$\mathrm{Hom}_{\mathcal{A}}(Q,{}_-)$ to this short exact sequence produces the usual long exact sequences in homology
for the categories $\mathcal{B}$ and $\mathcal{A}$, respectively.
Moreover, the homomor\-phisms~$\iota^d_{Q,{}_-}$ give rise to a~homomorphism between these long exact sequences.
A~similar statement is true for $\mathrm{Hom}_{\mathcal{B}}({}_-,Q)$ and $\mathrm{Hom}_{\mathcal{A}}({}_-,Q)$.

\subsection{Checking extension fullness}%\label{s1.3}

In this section we formulate and prove three propositions which will be useful for our study of extension fullness later
in the paper.
The following statement is a~modif\/ication of~\cite[Lemma~4]{CoMa} which also allows for a~somewhat stronger formulation.

\begin{proposition}
\label{prop1}
Let $\mathcal{A}$ be an Abelian category and $\mathcal{B}$ a~Serre subcategory of $\mathcal{A}$.
Assume that $\mathcal{B}$ contains a~full subcategory $\mathcal{B}_0$ with the following properties:
\begin{enumerate}[$($a$)$]\itemsep=0pt
\item
\label{prop1.1}
Every object in $\mathcal{B}$ is a~quotient of an object in $\mathcal{B}_0$.
\item
\label{prop1.2}
The homomorphisms $\iota^d_{M,N}$ are bijective for all $d\in\mathbb{Z}_{\geq0}$, $M\in\mathcal{B}_0$ and
$N\in\mathcal{B}$.
\end{enumerate}
Then $\mathcal{B}$ is {\em extension full} in $\mathcal{A}$.
\end{proposition}

\begin{proof}
Our proof is similar to that of~\cite[Lemma~4]{CoMa}.
We prove the statement by induction on~$d$.
Since $\mathcal{B}$ is assumed to be a~Serre subcategory of $\mathcal{A}$, it is clear that $\iota^0_{Q,N}$ and
$\iota^1_{Q,N}$ are isomorphisms for all $Q,N\in\mathcal{B}$.

To prove the induction step, for $Q\in\mathcal{B}$ consider a~short exact sequence
\begin{gather*}
0\to K\to M\to Q\to 0
\end{gather*}
with $M\in\mathcal{B}_0$ and $K\in\mathcal{B}$, which exists by condition~\eqref{prop1.1}.
Applying $\mathrm{Hom}_{\mathcal{B}}({}_-,N)$ and $\mathrm{Hom}_{\mathcal{A}}({}_-,N)$ gives for each~$d$ the following
commutative diagram with exact rows:
\begin{gather*}
\xymatrix{
\mathrm{Ext}^{d-1}_{\mathcal{B}}(M,N)\ar[r]\ar[d]_{\iota^{d-1}_{M,N}}
&\mathrm{Ext}^{d-1}_{\mathcal{B}}(K,N)\ar[r]\ar[d]_{\iota^{d-1}_{K,N}}
&\mathrm{Ext}^d_{\mathcal{B}}(Q,N)\ar[r]\ar[d]_{\iota^d_{Q,N}}
& \mathrm{Ext}^d_{\mathcal{B}}(M,N)\ar[d]_{\iota^d_{M,N}}\\
\mathrm{Ext}^{d-1}_{\mathcal{A}}(M,N)\ar[r]&\mathrm{Ext}^{d-1}_{\mathcal{A}}(K,N)\ar[r]
&\mathrm{Ext}^d_{\mathcal{A}}(Q,N)\ar[r]& \mathrm{Ext}^d_{\mathcal{A}}(M,N)
}
\end{gather*}
Now, $\iota^{d-1}_{M,N}$ and $\iota^d_{M,N}$ are isomorphisms by condition~\eqref{prop1.2} and from the induction step
we have that $\iota^{d-1}_{K,N}$ is a~monomorphism.
The injective four lemma (by which we mean the statement of~\cite[Lemma~I.3.3(i)]{Mc}) therefore implies that
$\iota^d_{Q,N}$ is a~monomorphism.

It remains to show that $\iota^d_{Q,N}$ is an epimorphism.
For this we consider the following commutative diagram with exact rows:
\begin{gather*}
\xymatrix{
\mathrm{Ext}^{d-1}_{\mathcal{B}}(K,N)\ar[r]\ar[d]_{\iota^{d-1}_{K,N}}
&\mathrm{Ext}^{d}_{\mathcal{B}}(Q,N)\ar[r]\ar[d]_{\iota^{d}_{Q,N}}
&\mathrm{Ext}^d_{\mathcal{B}}(M,N)\ar[r]\ar[d]_{\iota^d_{M,N}}
& \mathrm{Ext}^d_{\mathcal{B}}(K,N)\ar[d]_{\iota^d_{K,N}}\\
\mathrm{Ext}^{d-1}_{\mathcal{A}}(K,N)\ar[r]&\mathrm{Ext}^{d}_{\mathcal{A}}(Q,N)\ar[r]
&\mathrm{Ext}^d_{\mathcal{A}}(M,N)\ar[r]& \mathrm{Ext}^d_{\mathcal{A}}(K,N)
}
\end{gather*}
As $\iota^{d-1}_{K,N}$ is a~bijection by the induction step, $\iota^d_{M,N}$ is a~bijection by condition~\eqref{prop1.2}
and $\iota^d_{K,N}$ is a~monomorphism by the previous paragraph, the surjective four lemma (by which we mean the
statement of~\cite[Lemma~I.3.3(ii)]{Mc}) implies that $\iota^d_{Q,N}$ is an epimorphism.
This completes the proof.
\end{proof}
\setcounter{section}{2}
\setcounter{subsection}{3}
\setcounter{subsubsection}{0}
\setcounter{innercustomthm}{0}
\setcounter{equation}{2}
\setcounter{proposition}{3}
% Original source lines 437--769
\subsection{Adjunction lemma}%\label{s1.4}

The following statement is standard when dealing with categories with enough projective or injective objects.
We failed to f\/ind it in the literature in the generality we need, so we provide a~proof without the use of projective or
injective objects.

\begin{proposition}[adjunction lemma]
\label{adjlem}
Let $\mathcal{A}$ and $\mathcal{B}$ be two Abelian categories and $(\mathrm{F},\mathrm{G})$ an adjoint pair of exact
functors $\mathrm{F}:\mathcal{A}\to \mathcal{B}$ and $\mathrm{G}:\mathcal{B}\to \mathcal{A}$.
Then for every $d\in\mathbb{Z}_{\geq 0}$, $N\in \mathcal{A}$ and $M\in \mathcal{B}$ there are isomorphisms
\begin{gather*}
\mathrm{Ext}^d_{\mathcal{B}}(\mathrm{F}(N),M)\cong \mathrm{Ext}^d_{\mathcal{A}}(N,\mathrm{G}(M))
\end{gather*}
natural in both~$N$ and~$M$.
\end{proposition}

\begin{proof}
Applying $\mathrm{F}$ to an exact sequence
\begin{gather*}
\xymatrix@C=7mm{
0\ar[r]&\mathrm{G}(M)\ar[r]&X_1\ar[r]&X_2\ar[r]&\cdots\ar[r]&X_{d-1}\ar[r]&X_d\ar[r]&N\ar[r]&0
}
\end{gather*}
in $\mathcal{A}$, gives an exact sequence
\begin{gather*}
\xymatrix@C=5mm{
0\ar[r]&\mathrm{F}\mathrm{G}(M)\ar[r]&\mathrm{F}(X_1)\ar[r]&
\mathrm{F}(X_2)\ar[r]&\cdots\ar[r]&\mathrm{F}(X_{d-1})\ar[r]&\mathrm{F}(X_d)\ar[r]&\mathrm{F}(N)\ar[r]&0
}
\end{gather*}
in $\mathcal{B}$.
Denote by $\mathrm{K}$ the kernel of the adjunction natural transformation $\mathrm{F}\mathrm{G}\to
\mathrm{Id}_{\mathcal{B}}$.
Then we have the following commutative diagram
\begin{gather}
\label{eq11}
\begin{split}
& \xymatrix@C=5mm{
0\ar[r]&\mathrm{K}(M)\ar[r]\ar@{^{(}->}[d]&\mathrm{K}(M)\ar[r]\ar@{^{(}->}[d]&
0\ar[d]\ar[r]&\cdots\ar[r]&0\ar[d]\ar[r]&0\ar[d]\ar[r]&0\ar[d]\ar[r]&0\\
0\ar[r]&\mathrm{F}\mathrm{G}(M)\ar[r]\ar@{->>}[d]&\mathrm{F}(X_1)\ar[r]\ar@{->>}[d]&
\mathrm{F}(X_2)\ar[r]\ar@{=}[d]&\cdots\ar[r]&\mathrm{F}(X_{d-1})\ar[r]\ar@{=}[d]&
\mathrm{F}(X_d)\ar[r]\ar@{=}[d]&\mathrm{F}(N)\ar[r]\ar@{=}[d]&0\\
0\ar[r]&M'\ar[r]\ar@{^{(}->}[d]&X'\ar[r]\ar@{^{(}->}[d]&
\mathrm{F}(X_2)\ar[r]\ar@{=}[d]&\cdots\ar[r]&\mathrm{F}(X_{d-1})\ar[r]\ar@{=}[d]&
\mathrm{F}(X_d)\ar[r]\ar@{=}[d]&\mathrm{F}(N)\ar[r]\ar@{=}[d]&0\\
0\ar[r]&M\ar[r]&X''\ar[r]&
\mathrm{F}(X_2)\ar[r]&\cdots\ar[r]&\mathrm{F}(X_{d-1})\ar[r]&\mathrm{F}(X_d)\ar[r]&\mathrm{F}(N)\ar[r]&0
}
\end{split}
\end{gather}
with exact rows.
Here the homomorphism from the f\/irst to the second row is given by the natural embedding $\mathrm{K}\hookrightarrow
\mathrm{F}\mathrm{G}$ and the third row is just the corresponding cokernel with the morphism from the second to the
third row being the canonical projection.
In particular, $M':=\mathrm{F}\mathrm{G}(M)/\mathrm{K}(M)$.
The homomorphism from $M'$ to~$M$ is the natural inclusion (coming from the def\/inition of $\mathrm{K}$) and, f\/inally,
$X''$ is def\/ined as the push-out and the map to $\mathrm{F}(X_2)$ is given by the universal property of push-outs.
Functoriality of the construction yields a~group homomorphism
\begin{gather*}
\Phi: \ \mathrm{Ext}^d_{\mathcal{A}}(N,\mathrm{G}(M))\to \mathrm{Ext}^d_{\mathcal{B}}(\mathrm{F}(N),M).
\end{gather*}

Applying $\mathrm{G}$ to an exact sequence
\begin{gather*}
\xymatrix@C=7mm{
0\ar[r]&M\ar[r]&Y_1\ar[r]&Y_2\ar[r]&\cdots\ar[r]&Y_{d-1}\ar[r]&Y_d\ar[r]&\mathrm{F}(N)\ar[r]&0
}
\end{gather*}
in $\mathcal{B}$, gives an exact sequence
\begin{gather*}
\xymatrix@C=5mm{
0\ar[r]&\mathrm{G}(M)\ar[r]&\mathrm{G}(Y_1)\ar[r]&\mathrm{G}(Y_2)\ar[r]&
\cdots\ar[r]&\mathrm{G}(Y_{d-1})\ar[r]&\mathrm{G}(Y_d)\ar[r]&\mathrm{G}\mathrm{F}(N)\ar[r]&0
}
\end{gather*}
in $\mathcal{A}$.
Denote by $\mathrm{C}$ the cokernel of the adjunction natural transformation $\mathrm{Id}_{\mathcal{A}}\to
\mathrm{G}\mathrm{F}$.
Then we have the following commutative diagram
\begin{gather*}%\label{eq12}
\xymatrix@C=5mm{
0\ar[r]&\mathrm{G}(M)\ar[r]\ar@{=}[d]&\mathrm{G}(Y_1)\ar[r]\ar@{=}[d]&\mathrm{G}(Y_2)\ar[r]\ar@{=}[d]&
\cdots\ar[r]&\mathrm{G}(Y_{d-1})\ar[r]\ar@{=}[d]&
Y''\ar[r]\ar@{->>}[d]&N\ar[r]\ar@{->>}[d]&0\\
0\ar[r]&\mathrm{G}(M)\ar[r]\ar@{=}[d]&\mathrm{G}(Y_1)\ar[r]\ar@{=}[d]&\mathrm{G}(Y_2)\ar[r]\ar@{=}[d]&
\cdots\ar[r]&\mathrm{G}(Y_{d-1})\ar[r]\ar@{=}[d]&
Y'\ar[r]\ar@{^{(}->}[d]&N'\ar[r]\ar@{^{(}->}[d]&0\\
0\ar[r]&\mathrm{G}(M)\ar[r]\ar[d]&\mathrm{G}(Y_1)\ar[r]\ar[d]&\mathrm{G}(Y_2)\ar[r]\ar[d]&
\cdots\ar[r]&\mathrm{G}(Y_{d-1})\ar[r]\ar[d]&
\mathrm{G}(Y_d)\ar[r]\ar@{->>}[d]&\mathrm{G}\mathrm{F}(N)\ar[r]\ar@{->>}[d]&0\\
0\ar[r]&0\ar[r]&0\ar[r]&0\ar[r]&
\cdots\ar[r]&0\ar[r]&\mathrm{C}(N)\ar[r]&\mathrm{C}(N)\ar[r]&0
}
\end{gather*}
with exact rows.
Here the homomorphism from the third to the last row is given by the natural projection of
$\mathrm{G}\mathrm{F}\twoheadrightarrow \mathrm{C}$ and the second row is just the corresponding kernel with the
morphism from the second to the third row being the canonical injection.
In particular, $\mathrm{C}(N):=\mathrm{G}\mathrm{F}(M)/N'$.
The homomorphism from~$N$ to $N'$ is the natural surjection and, f\/inally, $Y''$ is def\/ined as the pullback and the map
from $\mathrm{G}(Y_{d-1})$ is given by the universal property of pullbacks.
Functoriality of the construction yields a~group homomorphism
\begin{gather*}
\Psi: \ \mathrm{Ext}^d_{\mathcal{B}}(\mathrm{F}(N),M)\to \mathrm{Ext}^d_{\mathcal{A}}(N,\mathrm{G}(M)).
\end{gather*}

Finally, we claim that~$\Phi$ and~$\Psi$ are inverses of each other.
Consider the following commutative diagram:
\begin{gather}
\label{eq14}
\begin{split}
& \xymatrix@C=3mm{
0\ar[r]&\mathrm{G}(M)\ar[r]\ar[d]&X_1\ar[r]\ar[d]&X_2\ar[r]\ar[d]&\cdots\ar[r]
&X_{d-1}\ar[r]\ar[d]&X_d\ar[r]\ar[d]&N\ar[r]\ar[d]&0\\
0\ar[r]&\mathrm{G}\mathrm{F}\mathrm{G}(M)\ar[r]\ar@{->>}[d]&\mathrm{G}\mathrm{F}(X_1)\ar[r]\ar@{->>}[d]&
\mathrm{G}\mathrm{F}(X_2)\ar[r]\ar@{=}[d]&\cdots\ar[r]&\mathrm{G}\mathrm{F}(X_{d-1})\ar[r]\ar@{=}[d]&
\mathrm{G}\mathrm{F}(X_d)\ar[r]\ar@{=}[d]&\mathrm{G}\mathrm{F}(N)\ar[r]\ar@{=}[d]&0\\
0\ar[r]&\mathrm{G}(M)\ar[r]&\mathrm{G}(X'')\ar[r]&
\mathrm{G}\mathrm{F}(X_2)\ar[r]&\cdots\ar[r]&\mathrm{G}\mathrm{F}(X_{d-1})\ar[r]&
\mathrm{G}\mathrm{F}(X_d)\ar[r]&\mathrm{G}\mathrm{F}(N)\ar[r]&0\\
0\ar[r]&\mathrm{G}(M)\ar[r]\ar@{=}[u]&\mathrm{G}(X'')\ar[r]\ar@{=}[u]&
\mathrm{G}\mathrm{F}(X_2)\ar[r]\ar@{=}[u]&\cdots\ar[r]&\mathrm{G}\mathrm{F}(X_{d-1})\ar[r]\ar@{=}[u]&
Y''\ar[r]\ar[u]&N\ar[r]\ar[u]&0
}
\end{split}
\end{gather}
Here the second row is obtained from the f\/irst one by applying $\mathrm{G}\mathrm{F}$ and the homomorphism from the
f\/irst to the second row is given by adjunction $\mathrm{Id}_{\mathcal{A}}\to\mathrm{G}\mathrm{F}$.
The second and the third rows and the homomorphism between them are given by applying the exact functor $\mathrm{G}$ to
the two middle rows of~\eqref{eq11}.
Note that, by construction, application of $\mathrm{G}$ identif\/ies the two last rows of~\eqref{eq11}, that is
$\mathrm{G}(M')\cong \mathrm{G}(M)$ and $\mathrm{G}(X')\cong \mathrm{G}(X'')$ (this follows from surjectivity of the
natural transformation $\mathrm{G}\mathrm{F}\mathrm{G}\to \mathrm{G}$ given by adjunction).
Consequently, from the adjunction identities we have that the composition of the maps in the f\/irst column is the
identity on $\mathrm{G}(M)$.
Finally, the last row and the homomorphism to the third row are given by the def\/inition of~$\Psi$.
In particular, from the construction we have that the image of~$N$ in the second row (coming from the f\/irst row) and in
the third row (coming from the last row) coincide.
Hence diagram~\eqref{eq14} gives rise to a~diagram
\begin{gather*}
\xymatrix@C=3mm{
0\ar[r]&\mathrm{G}(M)\ar[r]\ar@{=}[d]&X_1\ar[r]\ar[d]&X_2\ar[r]\ar[d]&\cdots\ar[r]
&X_{d-1}\ar[r]\ar[d]&X_d\ar[r]\ar[d]&N\ar[r]\ar[d]&0\\
0\ar[r]&\mathrm{G}(M)\ar[r]&\mathrm{G}(X'')\ar[r]&
\mathrm{G}\mathrm{F}(X_2)\ar[r]&\cdots\ar[r]&\mathrm{G}\mathrm{F}(X_{d-1})\ar[r]&
Q'\ar[r]&N'\ar[r]&0\\
0\ar[r]&\mathrm{G}(M)\ar[r]\ar@{=}[u]&\mathrm{G}(X'')\ar[r]\ar@{=}[u]&
\mathrm{G}\mathrm{F}(X_2)\ar[r]\ar@{=}[u]&\cdots\ar[r]&\mathrm{G}\mathrm{F}(X_{d-1})\ar[r]\ar@{=}[u]&
Y''\ar[r]\ar[u]&N\ar[r]\ar[u]&0
}
\end{gather*}
with exact rows, where $N'$ is the image of~$N$ in $\mathrm{G}\mathrm{F}(N)$ and $Q'$ is the full preimage of $N'$.
Pulling back along the epimorphism $N\twoheadrightarrow N'$ gives a~commutative diagram
\begin{gather*}
\xymatrix@C=3mm{
0\ar[r]&\mathrm{G}(M)\ar[r]\ar@{=}[d]&X_1\ar[r]\ar[d]&X_2\ar[r]\ar[d]&\cdots\ar[r]
&X_{d-1}\ar[r]\ar[d]&X_d\ar[r]\ar[d]&N\ar[r]\ar@{=}[d]&0\\
0\ar[r]&\mathrm{G}(M)\ar[r]&\mathrm{G}(X'')\ar[r]&
\mathrm{G}\mathrm{F}(X_2)\ar[r]&\cdots\ar[r]&\mathrm{G}\mathrm{F}(X_{d-1})\ar[r]&
Q\ar[r]&N\ar[r]&0\\
0\ar[r]&\mathrm{G}(M)\ar[r]\ar@{=}[u]&\mathrm{G}(X'')\ar[r]\ar@{=}[u]&
\mathrm{G}\mathrm{F}(X_2)\ar[r]\ar@{=}[u]&\cdots\ar[r]&\mathrm{G}\mathrm{F}(X_{d-1})\ar[r]\ar@{=}[u]&
Y''\ar[r]\ar[u]&N\ar[r]\ar@{=}[u]&0
}
\end{gather*}
with exact rows.
The last diagram shows that the extensions given by the f\/irst and the last rows coincide, which proves that $\Psi\Phi$
is the identity map.

The claim that $\Phi\Psi$ is the identity map is proved similarly.
This completes the proof.
\end{proof}

\section{Gelfand--Zeitlin modules}\label{s2}

\textbf{Notation.} For a~Lie algebra $\mathfrak{a}$ we denote by $U(\mathfrak{a})$ its universal enveloping algebra and
by $Z(\mathfrak{a})$ the center of $U(\mathfrak{a})$.

\subsection[Gelfand--Zeitlin subalgebra of $\mathfrak{gl}_n$]{Gelfand--Zeitlin subalgebra of $\boldsymbol{\mathfrak{gl}_n}$}%\label{s2.1}

For $k\in\mathbb{Z}_{>0}$ denote by $\mathfrak{g}_k$ the Lie algebra $\mathfrak{gl}_k(\mathbb{C})$.
Set $U_k=U(\mathfrak{g}_k)$ and $Z_k=Z(\mathfrak{g}_k)$.
We consider the usual chain
\begin{gather*}
\mathfrak{g}_1\subset \mathfrak{g}_2\subset \mathfrak{g}_3\subset \dots \subset \mathfrak{g}_n
\end{gather*}
of ``left upper corner'' embeddings of Lie algebras as depicted on the following picture:

\begin{center}
\begin{picture}(143.00,143.00)
\drawline(20,20)(140,20)
\drawline(20,20)(20,140)
\drawline(140,140)(20,140)
\drawline(140,140)(140,20)
\drawline(20,120)(40,120)
\drawline(40,140)(40,120)
\drawline(20,100)(60,100)
\drawline(60,140)(60,100)
\drawline(20,80)(80,80)
\drawline(80,140)(80,80)
\drawline(20,40)(120,40)
\drawline(120,140)(120,40)
\put(25,130){$\mathfrak{g}_1$}
\put(45,110){$\mathfrak{g}_2$}
\put(65,90){$\mathfrak{g}_3$}
\put(85,70){$\cdot$}
\put(87.5,67.5){$\cdot$}
\put(90,65){$\cdot$}
\put(100,50){$\mathfrak{g}_{n\text{-}1}$}
\put(125,30){$\mathfrak{g}_n$}
\end{picture}
\end{center}

\vspace{-6mm}

This gives rise to the chain
\begin{gather*}
U_1\subset U_2\subset U_3\subset \dots \subset U_n
\end{gather*}
of embeddings of associative algebras.
The subalgebra~$\Gamma$ of $U:=U_n$ generated by all centers~$Z_k$, where $k=1,2,\dots,n$, is called the {\em
Gelfand--Zeitlin} subalgebra.
We set $\mathfrak{g}:=\mathfrak{g}_n$.

The algebra~$\Gamma$ is obviously commutative, moreover,~$\Gamma$ is a~polynomial algebra in $\frac{n(n+1)}{2}$
variables (these can be taken to be generators of $Z_k$ for $k=1,2,\dots,n$), see~\cite{GZ} or~\cite[Chapter~X]{Zh}.
Considered as a~subalgebra of~$U$,~$\Gamma$ is a~{\em Harish-Chandra subalgebra} in the sense of~\cite[Section~1]{DFO},
which means that every f\/initely generated~$\Gamma$-subbimodule of~$U$ is already f\/initely generated both as a~left and
as a~right~$\Gamma$-module, see~\cite[Theorem~24]{DFO}.
Furthermore,~$U$ is free both as a~left and as a~right~$\Gamma$-module, see~\cite{Ov}.
Consequently, the usual induction and coinduction functors
\begin{gather*}
\mathrm{Ind}_{\Gamma}^U:=U\bigotimes_{\Gamma}{}_-
\qquad
\text{and}
\qquad
\mathrm{Coind}_{\Gamma}^U:=\mathrm{Hom}_{\Gamma}(U,{}_-)
\end{gather*}
are exact.

\subsection{Gelfand--Zeitlin modules}%\label{s2.2}

The category $\mathscr{G}\mathscr{Z}$ of {\em Gelfand--Zeitlin}-modules for~$U$ is def\/ined as the full subcategory of
$U\dash{mod}$ (the category of f\/initely generated~$U$-modules), which consists of those $M\in U\dash{mod}$ on which the
action of~$\Gamma$ is locally f\/inite, that is $\dim(\Gamma v)<\infty$ for all $v\in M$.
The category $\mathscr{G}\mathscr{Z}$ is a~Serre subcategory of $U\dash{mod}$ as~$U$ is Noetherian.

Write $\mathrm{Specm}(\Gamma)$ for the set of maximal ideals in~$\Gamma$.
As~$\Gamma$ is a~polynomial algebra, we have the decomposition
\begin{gather*}
\Gamma\dash{fmod}=\bigoplus_{\mathbf{m}\in\mathrm{Specm}(\Gamma)} \Gamma\dash{fmod}_{\mathbf{m}},
\end{gather*}
where $\Gamma\dash{fmod}_{\mathbf{m}}$ denotes the full subcategory of $\Gamma\dash{fmod}$ consisting of all objects
annihilated by some power of $\mathbf{m}$.
From this decomposition, consider the functor
\begin{gather*}
\mathrm{Ind}_{\Gamma,\mathbf{m}}^U: \ \Gamma\dash{fmod}_{\mathbf{m}}\to \mathscr{G}\mathscr{Z}
\end{gather*}
def\/ined as the restriction of $\mathrm{Ind}_{\Gamma}^U$ to $\Gamma\dash{fmod}_{\mathbf{m}}$.

Note that modules in $\mathscr{G}\mathscr{Z}$ are usually inf\/inite-dimensional and hence the usual restriction
$\mathrm{Res}_{\Gamma}^U$ ends up in $\Gamma\dash{lfmod}$ and not in $\Gamma\dash{fmod}$.
However, from~\cite[Corollary~5.3(a)]{FuOv2} it follows that for any $M\in \mathscr{G}\mathscr{Z}$ and
$\mathbf{m}\in\mathrm{Specm}(\Gamma)$ the space
\begin{gather*}
M_{\mathbf{m}}:=\big\{v\in \mathrm{Res}_{\Gamma}^U M: \mathbf{m}^iv=0~\text{for some}~i\big\}
\end{gather*}
is f\/inite-dimensional.
This allows us to def\/ine the functor
\begin{gather*}
\mathrm{Res}_{\Gamma,\mathbf{m}}^U: \ \mathscr{G}\mathscr{Z}\to \Gamma\dash{fmod}_{\mathbf{m}},
\end{gather*}
which sends~$M$ to $M_{\mathbf{m}}$ and is def\/ined on morphisms by restriction.
Then the usual adjunction between induction and restriction implies that $\mathrm{Ind}_{\Gamma,\mathbf{m}}^U$ is left
adjoint to $\mathrm{Res}_{\Gamma,\mathbf{m}}^U$.
The functor $\mathrm{Res}_{\Gamma,\mathbf{m}}^U$ is, in turn, left adjoint to the restriction
\begin{gather*}
\mathrm{Coind}_{\Gamma,\mathbf{m}}^U: \ \Gamma\dash{fmod}_{\mathbf{m}}\to \mathscr{G}\mathscr{Z}
\end{gather*}
of the coinduction functor to $\Gamma\dash{fmod}_{\mathbf{m}}$.

\subsection{The main result}%\label{s2.3}

Our main result in this section is the following:

\begin{theorem}
\label{thm3}
The category $\mathscr{G}\mathscr{Z}$ is extension full in $U\dash{Mod}$.
\end{theorem}

\begin{proof}
To prove Theorem~\ref{thm3}, we would like to apply Proposition~\ref{prop1} for $\mathcal{A}=U\dash{Mod}$,
$\mathcal{B}=\mathscr{G}\mathscr{Z}$ and~$\mathcal{B}_0$ being the full subcategory of $\mathcal{B}$ consisting of
all~$U$-modules isomorphic to $\mathrm{Ind}_{\Gamma}^U N$ for a~f\/inite-dimensional~$\Gamma$-module~$N$.
Taking the assumptions of Proposition~\ref{prop1} into account, Theorem~\ref{thm3} reduces to the following lemma:

\begin{lemma}
%\label{lem4}
Let $Q\in\mathscr{G}\mathscr{Z}$,~$N$ be a~finite-dimensional~$\Gamma$-module and $M=\mathrm{Ind}_{\Gamma}^U N$.
Then $\iota^d_{M,Q}$ is an isomorphism for every $d\in\mathbb{Z}_{\geq 0}$.
\end{lemma}

\begin{proof}
By additivity, we may assume $N\in \Gamma\dash{fmod}_{\mathbf{m}}$ for some $\mathbf{m}\in\mathrm{Specm}(\Gamma)$.
The image of the functor $\mathrm{Ind}_{\Gamma,\mathbf{m}}^U:\Gamma\dash{fmod}_{\mathbf{m}}\to \mathcal{A}$ belongs to
$\mathcal{B}$.
We show that for every $d\in\mathbb{Z}_{\geq 0}$, any $N\in \Gamma\dash{fmod}_{\mathbf{m}}$ and any
$Q\in\mathscr{G}\mathscr{Z}$ we have the isomorphisms
\begin{gather}
\label{eq1}
\mathrm{Ext}^d_{\mathcal{A}}\big(\mathrm{Ind}_{\Gamma,\mathbf{m}}^U N, Q\big)  \cong
\mathrm{Ext}^d_{\Gamma\dash{Mod}}\big(N,\mathrm{Res}_{\Gamma}^U Q\big),
\\
\label{eq2}
\mathrm{Ext}^d_{\Gamma\dash{Mod}}\big(N,\mathrm{Res}_{\Gamma}^U Q\big)  \cong
\mathrm{Ext}^d_{\Gamma\dash{mod}_{\mathbf{m}}}\big(N,\mathrm{Res}_{\Gamma,\mathbf{m}}^U Q\big),
\\
\label{eq3}
\mathrm{Ext}^d_{\Gamma\dash{mod}_{\mathbf{m}}}\big(N,\mathrm{Res}_{\Gamma,\mathbf{m}}^U Q\big)  \cong
\mathrm{Ext}^d_{\mathcal{B}}\big(\mathrm{Ind}_{\Gamma,\mathbf{m}}^U N, Q\big).
\end{gather}
Since, by construction, these three isomorphisms together with the morphism $\iota^d_{M,Q}$ will yield a~commutative
square, we get that $\iota^d_{M,Q}$ is an isomorphism.

Isomorphism~\eqref{eq2} follows from~\cite[Lemma~17]{CoMa}.
Isomorphisms~\eqref{eq1} and~\eqref{eq3} follow from adjunction lemma (Proposition~\ref{adjlem}).
Note that adjunction lemma applies here since~$U$ is free over~$\Gamma$.
\end{proof}

Theorem~\ref{thm3} now follows.
\end{proof}
\begin{thebibliography}{99}
\bibitem[7]{CoMa}
Coulembier K., Mazorchuk V., Some homological properties of the category
  {${\mathcal O}$}.~{III}, \href{http://arxiv.org/abs/1404.3401}{arXiv:1404.3401}.

\bibitem[12]{DFO}
Drozd Yu.A., Futorny V.M., Ovsienko S.A., Harish-{C}handra subalgebras and
  {G}el'fand--{Z}etlin modules, in Finite-Dimensional Algebras and Related
  Topics ({O}ttawa, {ON}, 1992), \href{http://dx.doi.org/10.1007/978-94-017-1556-0_5}{\textit{NATO Adv. Sci. Inst. Ser. C Math.
  Phys. Sci.}}, Vol.~424, Kluwer Acad. Publ., Dordrecht, 1994, 79--93.

\bibitem[21]{FuOv2}
Futorny V., Ovsienko S., Fibers of characters in {G}elfand--{T}setlin
  categories, \href{http://dx.doi.org/10.1090/S0002-9947-2014-05938-2}{\textit{Trans. Amer. Math. Soc.}} \textbf{366} (2014), 4173--4208,
  \href{http://arxiv.org/abs/math.RT/0610071}{math.RT/0610071}.

\bibitem[23]{GZ}
Gel'fand I.M., Tsetlin M.L., Finite-dimensional representations of the group of
  unimodular matrices, \textit{Dokl. Akad. Nauk USSR} \textbf{71} (1950),
  825--828.

\bibitem[31]{Mc}
Mac~Lane S., Homology, \textit{Classics in Mathematics}, Springer-Verlag, Berlin, 1995.

\bibitem[44]{Ov}
Ovsienko S., Finiteness statements for {G}elfand--{Z}etlin modules, in Third
  {I}nternational {A}lgebraic {C}onference in the {U}kraine, Inst. of Math.,
  Kiev, 2002, 323--338.

\bibitem[49]{We}
Weibel C.A., An introduction to homological algebra, \href{http://dx.doi.org/10.1017/CBO9781139644136}{\textit{Cambridge Studies
  in Advanced Mathematics}}, Vol.~38, Cambridge University Press, Cambridge,
  1994.

\bibitem[50]{Zh}
Zhelobenko D.P., Compact Lie groups and their representations, Nauka, Moscow,
  1970.

\bibitem[51]{Zi}
Zimmermann A., Representation theory. A~homological algebra point of view,
  \href{http://dx.doi.org/10.1007/978-3-319-07968-4}{\textit{Algebra and Applications}}, Vol.~19, Springer International
  Publishing, Switzerland, 2014.

\end{thebibliography}
\end{document}
